% TOP_PROBLEM: {"id": 3172, "title": "Hamiltonicity of Cayley graphs", "category_id": 3, "category": {"id": 3, "name": "graph_theory", "display_name": "Graph Theory", "description": "Problems involving graphs, networks, and their properties.", "slug": "graph-theory", "order_index": 3, "created_at": "2026-07-31T15:26:25.671Z"}, "status": "open", "problem_number": "OPG-2095", "set_id": 12, "statusQualification": "The terse source is interpreted in its standard finite connected undirected form, with a generating inverse-closed set and group order at least three."}
% FIRST_POSTED: 2026-10-01
% ENTRY_KIND: statement_only
% UnsolvedMath ID 3172; source catalogue code OPG-2095
% !TeX program = lualatex
% Definition and sources only; this file is not an LLM solution attempt.
\documentclass[11pt,a4paper]{article}
\usepackage[margin=25mm]{geometry}
\usepackage{fontspec}
\setmainfont{lmroman10-regular.otf}[BoldFont=lmroman10-bold.otf,
 ItalicFont=lmroman10-italic.otf,BoldItalicFont=lmroman10-bolditalic.otf]
\usepackage{amsmath,amssymb,mathtools}
\usepackage{unicode-math}
\setmathfont{latinmodern-math.otf}
\AtBeginDocument{\let\setminus\smallsetminus}
\usepackage{xurl,hyperref}
\hypersetup{colorlinks=true,urlcolor=blue}
\setcounter{secnumdepth}{0}
\setlength{\parindent}{0pt}
\setlength{\parskip}{5pt}
\setlength{\emergencystretch}{3em}
\begin{document}
\section*{Hamiltonicity of Cayley graphs}\label{3172}
{\small UnsolvedMath ID 3172 · OPG-2095 · Graph Theory · OpenGarden}
\subsection{Definitions and mathematical statement}
Let $\Gamma$ be a finite group with identity $1$ and order
$N\geq3$. Choose an inverse-closed generating set
$S\subseteq\Gamma\setminus\{1\}$, so $S^{-1}=S$ and
$\langle S\rangle=\Gamma$. The undirected Cayley graph
$\operatorname{Cay}(\Gamma,S)$ has vertex set $\Gamma$ and
an edge between distinct $g,h$ precisely when
$g^{-1}h\in S$. Inverse closure makes it undirected,
excluding $1$ removes loops, and the generating condition
makes the graph connected. Repeated descriptions of an edge
give only one edge.

The Hamiltonicity question is whether every such graph
contains a Hamilton cycle. Explicitly, for every pair
$(\Gamma,S)$ there should exist an enumeration
\[
 g_1,g_2,\ldots,g_N
\]
of the group elements such that $g_i^{-1}g_{i+1}\in S$
for all indices modulo $N$.

The enumeration must use every group element exactly once
before returning to its start. The generating set is part
of the input and cannot be replaced by a more convenient
one. The group is not assumed abelian. This records the
standard finite connected interpretation of the source's
short question: an arbitrary disconnected Cayley graph
cannot have one spanning cycle, and orders one and two
do not contain a simple cycle under the convention here.
Directed Cayley graphs without inverse-closed sets constitute
a different problem and are not included.
\subsection{Short English statement}
Does every finite connected undirected Cayley graph with at least three vertices have a Hamilton cycle?
\subsection{Sources}
\begin{itemize}
\item[S1] ULAM.AI and the UnsolvedMath Contributors, supplied problems.json, record 3172 (OPG-2095), OpenGarden collection. Catalogue identity and source transcription; CC BY 4.0.
\url{https://huggingface.co/datasets/ulamai/UnsolvedMath}.
\item[S2] Open Problem Garden, Hamiltonicity of Cayley graphs. Original problem and discussion; the mathematical formulation below is expanded from the supplied source record.
\url{http://www.openproblemgarden.org/op/hamiltonicity_of_cayley_graphs}.
\end{itemize}
\end{document}
